\documentclass[10pt,a4paper]{amsart}
\usepackage{amsmath,amssymb,amsfonts,amsthm}
\usepackage{enumitem}
\usepackage[left=1.5cm, right=1.5cm]{geometry}

\newcommand{\R}{\mathbb{R}}
\renewcommand{\d}{\mathrm{d}}

\theoremstyle{definition}
\newtheorem{problema}{Problema}

\title{Introdução às Equações Diferenciais -- Lista 4}

\begin{document}
\maketitle

\noindent{\small Os problemas estão organizados por tema. A numeração é
contínua ao longo de toda a lista.}

\vspace{4mm}

%=====================================================================
\section*{1. Equações de primeira ordem}
%=====================================================================

\begin{problema}
Resolva a equação exata
\[
    (2xy+e^x)\,\d x+(x^2+\cos y)\,\d y=0,
    \qquad y(0)=0.
\]
Encontre a solução implicitamente e determine se ela define $y$ localmente
como função de $x$ em uma vizinhança de $x=0$.
\end{problema}

\begin{problema}
Resolva o problema de valor inicial
\[
    y'+\frac{2t}{1+t^2}\,y=\frac{1}{1+t^2},
    \qquad y(0)=2,
\]
e determine o intervalo maximal no qual a solução está definida.
\end{problema}

\begin{problema}
Resolva o problema de valor inicial
\[
    y'=x(1+y^2),
    \qquad y(0)=1,
\]
e determine o intervalo maximal no qual a solução está definida.
\end{problema}

\begin{problema}
Resolva
\[
    y'+y=e^{-t},
    \qquad y(0)=2,
\]
e descreva o comportamento da solução quando $t\to+\infty$.
\end{problema}

\begin{problema}
Resolva a equação exata
\[
    (y\cos x+2x)\,\d x+\sin x\,\d y=0,
    \qquad y(0)=1,
\]
exibindo a solução explicitamente sempre que possível.
\end{problema}

%=====================================================================
\section*{2. Existência, unicidade e iteração de Picard}
%=====================================================================

\begin{problema}
Descreva várias soluções distintas de
\[
    y'=2\sqrt{|y|}, \qquad y(0)=0,
\]
e explique precisamente qual hipótese do teorema de Picard--Lindelöf falha.
\end{problema}

\begin{problema}
Encontre todas as soluções de
\[
    y'=3y^{2/3}, \qquad y(0)=0,
\]
com $y(t)\ge 0$, e explique por que isto não contradiz o teorema de
Picard--Lindelöf.
\end{problema}

\begin{problema}
Considere a equação de segunda ordem
\[
    x''=\sin x+t\,x'.
\]
Reescreva-a como um sistema de primeira ordem e prove que, para todos
$a,b\in\R$, existe uma única solução local satisfazendo $x(0)=a$ e
$x'(0)=b$.
\end{problema}

\begin{problema}
Prove a fórmula de adição
\[
    f(x+y)=\frac{f(x)+f(y)}{1-f(x)f(y)}
\]
para a solução maximal de
\[
    f'=1+f^2, \qquad f(0)=0,
\]
válida sempre que ambos os lados estejam definidos e $f(x)f(y)\neq 1$. Em
seguida, resolva a equação explicitamente.
\end{problema}

\begin{problema}
Defina funções $u_0,u_1,u_2,\ldots$ em $[0,1/4]$ por
\[
    u_0(x)=0,
    \qquad
    u_{n+1}(x)=x+\int_0^x u_n(t)^2\,\d t.
\]
\begin{enumerate}[label=(\alph*)]
    \item Prove que $u_n$ converge uniformemente em $[0,1/4]$.
    \item Identifique a função limite determinando a equação diferencial que ela satisfaz.
\end{enumerate}
\end{problema}

\begin{problema}
Seja $n$ um inteiro maior que $1$. Existe uma função diferenciável em
$[0,\infty)$ cuja derivada é igual à sua $n$-ésima potência e cujo valor na
origem é positivo?
\end{problema}

\begin{problema}
Encontre todas as soluções diferenciáveis da equação
\[
    y'=\sqrt{y}, \qquad y(0)=0.
\]
\end{problema}

%=====================================================================
\section*{3. Equações autônomas e retrato de fase}
%=====================================================================

\begin{problema}
Para a equação autônoma
\[
    y'=y(y-1)(y+2),
\]
classifique os pontos estacionários e determine
$\lim_{t\to+\infty}y(t)$ para cada valor inicial $y(0)=y_0$.
\end{problema}

\begin{problema}
Considere a equação autônoma
\[
    y'=y^2(4-y^2).
\]
classifique todos os pontos de equilíbrio e
determine $\lim_{t\to+\infty}y(t)$ para todo valor inicial cuja solução existe para todo tempo positivo.
\end{problema}

%=====================================================================
\section*{4. Equações lineares de segunda ordem}
%=====================================================================

\begin{problema}
Resolva
\[
    y''-2y'+y=e^t(1+t),
    \qquad y(0)=0,\qquad y'(0)=1.
\]
\end{problema}

\begin{problema}
Sabendo que $y_1(t)=t$ é uma solução em $(0,\infty)$, encontre a solução
geral de
\[
    t^2y''-2ty'+2y=0.
\]
\end{problema}

\begin{problema}
Sejam $y_1,y_2$ soluções de
\[
    y''+\tan(t)\,y'+\cos(t)\,y=0
\]
em $(-\pi/2,\pi/2)$, cujo Wronskiano satisfaz $W(0)=2$. Encontre $W(t)$ e
decida se $y_1,y_2$ podem se tornar linearmente dependentes nesse intervalo.
\end{problema}

\begin{problema}
Resolva
\[
    y''+4y=3\sin(2t),
    \qquad y(0)=0,
    \qquad y'(0)=1.
\]
\end{problema}

\begin{problema}
Resolva
\[
    y''-3y'+2y=t\,e^t.
\]
\end{problema}

\begin{problema}
Sabendo que $y_1(x)=x$ é uma solução de
\[
    (1-x^2)y''-2xy'+2y=0
\]
em $(-1,1)$, encontre uma segunda solução linearmente independente nesse
intervalo.
\end{problema}

\begin{problema}
Sejam $p,q$ funções contínuas em um intervalo $I$, e seja $t_0\in I$. Prove
que, se $y$ resolve
\[
    y''+p(t)y'+q(t)y=0
\]
e satisfaz $y(t_0)=y'(t_0)=0$, então $y\equiv 0$ em $I$. Deduza que duas
soluções com o mesmo valor e a mesma derivada em $t_0$ coincidem.
\end{problema}

%=====================================================================
\section*{5. Sistemas lineares e exponencial de matrizes}
%=====================================================================

\begin{problema}
Sejam $A,B$ matrizes reais ou complexas. Prove que
\[
    e^{t(A+B)}=e^{tA}\cdot e^{tB}
\]
para todo $t\in\R$ se e somente se $AB=BA$.
\end{problema}

\begin{problema}
Calcule $e^{tA}$ e resolva $x'=Ax$ para
\[
    A=\begin{pmatrix}
        2 & 1 & 0\\
        0 & 2 & 0\\
        0 & 0 & -1
    \end{pmatrix}.
\]
\end{problema}

\begin{problema}
Encontre todas as funções diferenciáveis $f,g:\R\to\R$ satisfazendo
\[
    f(x+y)=f(x)f(y)-g(x)g(y),
    \qquad
    g(x+y)=f(x)g(y)+g(x)f(y),
\]
para todos $x,y\in\R$, juntamente com
\[
    f(0)=1,\quad g(0)=0,\quad f'(0)=0,\quad g'(0)=1.
\]
\end{problema}

\begin{problema}
Seja $A:[0,T]\to M_n(\R)$ contínua e seja $X:[0,T]\to M_n(\R)$ a solução de
\[
    X'(t)=A(t)X(t),
    \qquad X(0)=I.
\]
\begin{enumerate}
    \item[(a)] Prove a a identidade de Jacobi
$\tfrac{\d}{\d t}\det X=\det X\cdot\operatorname{tr}(X^{-1}X')$ enquanto $X$
for invertível.
\item[(b)] Prove a fórmula de Liouville
\[
    \det X(t)=\exp\left(\int_0^t \operatorname{tr}(A(s))\,\d s\right),
\]
e conclua, em particular, que $X(t)$ é invertível para todo $t$.
\end{enumerate}
\end{problema}

\begin{problema}
\begin{enumerate}[label=(\alph*)]
    \item Encontre uma base para o espaço das soluções reais da equação
    diferencial
    \[
        (*)\qquad \sum_{n=0}^{7}\frac{\d^n x}{\d t^n}=0.
    \]
    \item Encontre uma base para o subespaço das soluções reais de $(*)$ que
    satisfazem $\lim_{t\to+\infty}x(t)=0$.
\end{enumerate}
\end{problema}

\begin{problema}
Considere o sistema de equações diferenciais
\[
    \frac{\d x}{\d t}=y+tz,\qquad
    \frac{\d y}{\d t}=z+t^2x,\qquad
    \frac{\d z}{\d t}=x+e^ty.
\]
Prove que existe uma solução definida para todo $t\in[0,1]$ tal que
\[
    \begin{pmatrix}1&2&3\\4&5&6\\7&8&9\end{pmatrix}
    \begin{pmatrix}x(0)\\y(0)\\z(0)\end{pmatrix}
    =\begin{pmatrix}0\\0\\0\end{pmatrix}
    \qquad\text{e}\qquad
    \int_0^1\bigl(x(t)^2+y(t)^2+z(t)^2\bigr)\,\d t=1.
\]
\end{problema}

\begin{problema}
Considere matrizes $A,B\in\R^{2\times2}$ e a equação diferencial matricial
\[
    X'(t)=AX(t)+X(t)B,\qquad X(t)\in\R^{2\times2}.
\]
\begin{enumerate}[label=(\alph*)]
    \item Prove que $X(t)=e^{At}X(0)e^{Bt}$ é solução.
    \item Se
    $A=\begin{pmatrix}0&-\omega\\\omega&0\end{pmatrix}$,
    $B=\begin{pmatrix}\alpha&0\\0&\beta\end{pmatrix}$ e
    $X(0)=\begin{pmatrix}x_{11}&x_{12}\\x_{21}&x_{22}\end{pmatrix}$,
    escreva explicitamente a solução $X(t)$.
    \item Prove que, caso $B = - A,$ então para todo $k\in\mathbb{N}$, o traço
    $\operatorname{Tr}(X(t)^k)$ é constante ao longo do tempo.
\end{enumerate}
\end{problema}

%=====================================================================
\section*{6. Sistemas não-lineares e equações qualitativas}
%=====================================================================

\begin{problema}
Sejam $l_0,c,\alpha,g$ constantes positivas, e seja $x(t)$ a solução da
equação diferencial
\[
    \bigl((l_0+ct^\alpha)^2x'\bigr)'
    +g[l_0+ct^\alpha]\sin x=0,
    \qquad t\ge 0,\qquad -\tfrac{\pi}{2}<x<\tfrac{\pi}{2},
\]
satisfazendo as condições iniciais
\[
    x(t_0)=x_0,\qquad x'(t_0)=0.
\]
Esta é a equação do pêndulo matemático cujo comprimento varia de acordo com
a lei $l=l_0+ct^\alpha$. Prove que $x(t)$ está definida no intervalo
$[t_0,\infty)$; além disso, se $\alpha>2$, então, para todo $x_0\neq 0$,
existe $t_0$ tal que
\[
    \liminf_{t\to\infty}|x(t)|>0.
\]
\end{problema}

\begin{problema}
Prove que, se $a_i$ $(i=1,2,3,4)$ são constantes positivas, $a_2-a_4>2$ e
$a_1a_3-a_2>2$, então a solução $(x(t),y(t))$ do sistema
\[
    \dot{x}=a_1-a_2x+a_3xy,
    \qquad
    \dot{y}=a_4x-y-a_3xy
    \qquad (x,y\in\R),
\]
com as condições iniciais
\[
    x(0)=0,\qquad y(0)\ge a_1,
\]
é tal que a função $x(t)$ possui exatamente um máximo local estrito no
intervalo $[0,\infty)$.
\end{problema}

\begin{problema}
Considere um sistema autônomo
\[
    \frac{\d x_i}{\d t}=F_i(x_1,\dots,x_n),
\]
onde $F=(F_1,\dots,F_n):\R^n\to\R^n$ é um campo vetorial de classe $C^1$.
\begin{enumerate}[label=(\alph*)]
    \item Sejam $U$ e $V$ duas soluções em $a<t<b$. Supondo que
    $\langle DF(x)z,z\rangle\le 0$ para todos $x,z\in\R^n$, mostre que
    $|U(t)-V(t)|^2$ é uma função decrescente de $t$.
    \item Seja $W(t)$ uma solução definida para $t>0$. Supondo que
    $\langle DF(x)z,z\rangle\le -|z|^2$, mostre que existe $C\in\R^n$ tal que
    $\lim_{t\to\infty}W(t)=C$.
\end{enumerate}
\end{problema}

\begin{problema}
Considere a equação não-linear de Ermakov--Pinney
\[
    y''+y=\frac{1}{y^3},\qquad y(0)=1,\quad y'(0)=1.
\]
\begin{enumerate}[label=(\alph*)]
    \item A partir da equação, obtenha uma equação de primeira ordem e uma
    quantidade conservada ao longo do tempo.
    \item Use as condições iniciais para determinar o valor da quantidade
    conservada.
    \item Resolva a equação da quantidade conservada e encontre uma solução
    da equação inicial.
\end{enumerate}
\end{problema}

%=====================================================================
\section*{7. Comportamento assintótico e equações com retardo}
%=====================================================================

\begin{problema}
Seja $f(x)$ uma função real tal que
\[
    \lim_{x\to+\infty}\frac{f(x)}{e^x}=1
\]
e
\[
    |f''(x)|<c|f'(x)|
\]
para todo $x$ suficientemente grande. Prove que
\[
    \lim_{x\to+\infty}\frac{f'(x)}{e^x}=1.
\]
\end{problema}

\begin{problema}
Sejam $a(x)$ e $r(x)$ funções positivas e contínuas em $[0,\infty)$, e
suponha que
\[
    \liminf_{x\to\infty}\bigl(x-r(x)\bigr)>0.
\]
Suponha que $y(x)$ seja uma função contínua em toda a reta real,
diferenciável em $[0,\infty)$, e que satisfaça
\[
    y'(x)=a(x)y(x-r(x))
\]
em $[0,\infty)$. Prove que o limite
\[
    \lim_{x\to\infty} y(x)\exp\left\{-\int_0^x a(u)\,\d u\right\}
\]
existe e é finito.
\end{problema}

\begin{problema}
Seja $x:[0,\infty)\to\R$ uma função diferenciável que satisfaz
\[
    x'(t)
    =
    -2x(t)\sin^2 t
    +
    (2-|\cos t|+\cos t)
    \int_{t-1}^{t} x(s)\sin^2 s\,\d s
\]
em $[1,\infty)$. Prove que $x$ é limitada em $[0,\infty)$ e que
$\lim_{t\to\infty}x(t)=0$. A conclusão continua verdadeira para funções que
satisfazem
\[
    x'(t)
    =
    -2x(t)
    +
    (2-|\cos t|+\cos t)
    \int_{t-1}^{t}x(s)\,\d s\,?
\]
\end{problema}

\begin{problema}
Seja $y(t)$ uma solução real, definida para $0<t<\infty$, da equação
\[
    \frac{\d y}{\d t}=e^{-y}-e^{-3y}+e^{-5y}.
\]
Mostre que $y(t)\to+\infty$ quando $t\to+\infty$.
\end{problema}

%=====================================================================
\section*{8. Estimativas e equações integrais}
%=====================================================================

\begin{problema}
Seja $f:[0,\infty)\to\R$ uma função duas vezes diferenciável tal que
\[
    f(0)=0,\qquad f'(0)=1,
\]
e
\[
    f''(x)=\frac{1}{1+f(x)}
\]
para todo $x\in[0,1]$. Prove que
\[
    f(1)<\frac{3}{2}.
\]
\end{problema}

\begin{problema}
Seja $a>0$. Encontre todas as funções continuamente diferenciáveis
$\varphi:[a,\infty)\to\R$ tais que, para todo $x>a$,
\[
    \varphi(x)^2
    =
    \int_a^x \left(|\varphi(y)|^2+|\varphi'(y)|^2\right)\,\d y
    -(x-a)^3.
\]
\end{problema}

\begin{problema}
Seja $f_0:[0,1]\to\R$ uma função contínua. Para $n\ge 1$, defina
\[
    f_n(x)=\int_0^x f_{n-1}(t)\,\d t.
\]
Prove que a série $\sum_{n=1}^{\infty} f_n(x)$ converge para todo
$x\in[0,1]$, e encontre uma fórmula explícita para sua soma.
\end{problema}

\begin{problema}
Seja $f:[0,\infty)\to\R$ contínua e suponha que
\[
    f(x)=1+\int_0^x (x-t)\bigl(f(t)-2e^t\bigr)\,\d t
    \qquad(x\ge 0).
\]
Mostre que $f$ satisfaz um problema de valor inicial de segunda ordem e
resolva-o explicitamente.
\end{problema}

%=====================================================================
\section*{9. Transformada de Laplace}
%=====================================================================

\begin{problema}
\begin{enumerate}
    \item Encontre a transformada de Laplace das seguintes funções.
    \begin{enumerate}[label=(\alph*)]
        \item $f(t)=\sin(2t)\cos(2t)$.
        \item $f(t)=\cos^2(3t)$.
        \item $f(t)=t e^{2t}\sin(3t)$.
        \item $f(t)=(t+3)u_7(t)$.
        \item $f(t)=t^2u_3(t)$.
        \item $f(t)=\begin{cases}1, & 0\le t<2,\\ t^2-4t+4, & t\ge 2.\end{cases}$
        \item $f(t)=\begin{cases}t, & 0\le t<3,\\ 5, & t\ge 3.\end{cases}$
        \item $f(t)=\begin{cases}0, & t<\pi,\\ t-\pi, & \pi\le t<2\pi,\\ 0, & t\ge 2\pi.\end{cases}$
        \item $f(t)=\begin{cases}\cos(\pi t), & t<4,\\ 0, & t\ge 4.\end{cases}$
        \item $f(t)=\begin{cases}t, & 0\le t<1,\\ e^t, & t\ge 1.\end{cases}$
    \end{enumerate}
    \item Encontre a transformada inversa de Laplace.
    \begin{enumerate}[label=(\alph*)]
        \item $F(s)=\dfrac{1}{(s+1)(s^2-1)}$.
        \item $F(s)=\dfrac{2s+3}{s^2+4s+13}$.
        \item $F(s)=\dfrac{e^{-3s}}{s-2}$.
        \item $F(s)=\dfrac{1+e^{-2s}}{s^2+6}$.
    \end{enumerate}
\end{enumerate}
\end{problema}

\begin{problema}
Use a transformada de Laplace para resolver
\[
    y''+4y=u(t)-u(t-\pi),
    \qquad y(0)=0,\qquad y'(0)=0,
\]
onde $u$ denota a função degrau.
\end{problema}

\begin{problema}
Prove que, se uma função $f$ é de ordem exponencial $s_0$ e possui
transformada de Laplace $F(s)=\mathcal{L}[f(t)]$, então
$\mathcal{L}[t^n f(t)]$ existe para $s>s_0$ e
\[
    \mathcal{L}[t^n f(t)]=(-1)^n F^{(n)}(s),\qquad s>s_0,
\]
onde $F^{(n)}=\dfrac{\d^n}{\d s^n}F$.
\end{problema}

\begin{problema}
Resolva o problema de valor inicial
\[
    y''+ty'-2y=4,\qquad y(0)=0,\quad y'(0)=0.
\]
\end{problema}

\begin{problema}
Resolva a equação integro-diferencial
\[
    y''+2y'+y+\int_0^t y(\tau)e^{-(t-\tau)}\,\d\tau=0,
    \qquad y(0)=1,\quad y'(0)=-1.
\]
\end{problema}

\begin{problema}
Resolva
\[
    y''+4y'+3y=12e^{-2t},\qquad y(0)=0,\quad \int_0^{\infty}y(t)\,\d t=1.
\]
\end{problema}

%=====================================================================
\section*{10. Equações diferenciais e funções geradoras}
%=====================================================================

\begin{problema} Seja $S_n = \{1,\dots,n\}.$
Seja $I_n$ o número de involuções em $S_n$. Isto é, $I_n$ é o número de permutações de $S_n$ que são iguais \`a própria inversa. 

\begin{enumerate}
    \item[(a)] Prove que, se definirmos $I_0 = 1,$ então
    \[
 I_1=1,\qquad I_n=I_{n-1}+(n-1)I_{n-2}.
\]
\item[(b)] Defina 
$$F(t)=\sum_{n\ge 0}I_n\frac{t^n}{n!}.$$
Deduza uma equação diferencial satisfeita por $F$. 
\item[(c)] Identifique $F$ explicitamente.
\end{enumerate}
\end{problema}

\begin{problema}
Seja $B_0=1$ e
\[
    B_{n+1}=\sum_{k=0}^n \binom{n}{k}B_k.
\]
Defina $B(t)=\sum_{n\ge 0}B_n\frac{t^n}{n!}$. Deduza uma equação
diferencial para $B(t)$, resolva-a e, expandindo a resposta em série de
potências, deduza a fórmula de Dobinski
\[
    B_n=e^{-1}\sum_{m=0}^{\infty}\frac{m^n}{m!}.
\]
\end{problema}

\begin{problema}
Seja $D_n$ o número de permutações de $\{1,\ldots,n\}$ sem ponto fixo (os
chamados desarranjos).
\begin{enumerate}[label=(\alph*)]
    \item Prove que $D_0=1$, $D_1=0$ e $D_{n+1}=n(D_n+D_{n-1})$ para
    $n\ge 1$.
    \item Para $F(t)=\sum_{n\ge 0}D_n\frac{t^n}{n!}$, deduza uma equação
    diferencial de primeira ordem, resolva-a e conclua que
    \[
        D_n=n!\sum_{k=0}^n \frac{(-1)^k}{k!}.
    \]
\end{enumerate}
\end{problema}

\begin{problema}
Fixe um inteiro $m\ge 2$ e considere o grafo completo $K_m$ com vértices
$\{1,\ldots,m\}$. Um passeio de comprimento $n$ é uma sequência de $n$
arestas em que arestas consecutivas compartilham um vértice. Seja $a_n$ o
número de passeios de comprimento $n$ que começam e terminam no vértice $1$,
e $b_n$ o número de passeios de comprimento $n$ que começam no vértice $1$ e
terminam em um vértice fixo $j\neq 1$.
\begin{enumerate}[label=(\alph*)]
    \item Prove que
    \[
        a_{n+1}=(m-1)b_n,\qquad b_{n+1}=a_n+(m-2)b_n,
    \]
    com $a_0=1$ e $b_0=0$.
    \item Para as funções geradoras exponenciais $A(t)$ e $B(t)$ de $a_n$ e
    $b_n$, deduza o sistema de equações diferenciais correspondente.
    \item Resolva o sistema e obtenha fórmulas fechadas para $a_n$ e $b_n$.
\end{enumerate}
\end{problema}

\begin{problema}
Seja $C_n$ o número de maneiras de colocar parênteses corretamente em um
produto de $n+1$ fatores. Isto é, o número de maneiras de colocar parênteses num produto tal que todos os parênteses são abertos e depois fechados e tal que, para fechar um parêntese, é necessário fechar antes todos os outros neste contido. Note, por exemplo, que $C_0=C_1=1$ e $C_2=2$.
\begin{enumerate}[label=(\alph*)]
    \item Prove que $C_{n+1}=\sum_{k=0}^n C_kC_{n-k}$.
    \item Para $C(t)=\sum_{n\ge 0}C_nt^n$, mostre que $C(t)=1+tC(t)^2$. 
    \item Deduza que $C_n=\dfrac{1}{n+1}\dbinom{2n}{n}$.
\end{enumerate}
\end{problema}

\begin{problema}
Seja $a_0=0$, $a_1=1$ e $a_{n+2}=a_{n+1}+a_n$ $(n\ge 0)$.
\begin{enumerate}[label=(\alph*)]
    \item Encontre a função geradora ordinária $A(t)=\sum_{n\ge 0}a_nt^n$.
    \item Use a função geradora para provar a fórmula de Binet
    \[
        a_n=\frac{1}{\sqrt5}\left[
        \left(\frac{1+\sqrt5}{2}\right)^n
        -\left(\frac{1-\sqrt5}{2}\right)^n
        \right].
    \]
\end{enumerate}
\end{problema}

\begin{problema}
Seja $a_0=1$, $a_1=4$ e $a_{n+2}=5a_{n+1}-6a_n$ $(n\ge 0)$. Encontre a
função geradora ordinária $A(t)=\sum_{n\ge 0}a_nt^n$ e, a partir dela, uma
fórmula fechada para $a_n$.
\end{problema}

\begin{problema}
Seja $a_0=0$ e $a_{n+1}=2a_n+n$ $(n\ge 0)$. Encontre a função geradora
ordinária $A(t)=\sum_{n\ge 0}a_nt^n$ e, a partir dela, uma fórmula fechada
para $a_n$.
\end{problema}

\begin{problema}
Seja $a_0=1$ e
\[
    a_{n+1}=\sum_{k=0}^n \binom{n}{k}a_ka_{n-k}
    \qquad(n\ge 0).
\]
Para a função geradora exponencial $A(t)=\sum_{n\ge 0}a_n\frac{t^n}{n!}$,
deduza uma equação diferencial, resolva-a e determine $a_n$ explicitamente.
\end{problema}

%=====================================================================
\section*{11. Cálculo das variações}
%=====================================================================

\begin{problema}
Determine a função $y\in C^1([0,L])$ que minimiza o funcional
\[
    \int_0^L (y'(x))^2\,\d x,
\]
sujeita a $y(0)=y(L)=0$ e $\int_0^L y(x)\,\d x=A$, onde $A>0$ é uma
constante dada.
\end{problema}

\begin{problema}
Suponha que exista um minimizador positivo do funcional
\[
    T[y]=\int_{x_0}^{x_1}\frac{\sqrt{1+y'(x)^2}}{\sqrt{y(x)}}\,\d x,
    \qquad y>0.
\]
\begin{enumerate}
    \item[(a)]Mostre que todo extremal satisfaz
\[
    y\bigl(1+y'^2\bigr)=a
\]
para alguma constante $a>0$. 
 \item Integre esta equação e conclua que o minimizador é um arco de
    ciclóide
    \[
        x(\varphi)=a(\varphi-\sin\varphi),
        \qquad
        y(\varphi)=a(1-\cos\varphi),
    \]
    a curva descrita por um ponto de um círculo de raio $a$ que rola sobre o
    eixo $x$.
    \item Explique por que a constante $a$ e o intervalo de $\varphi$ ficam
    determinados pela exigência de que a curva passe por $(x_1,y_1)$.
\end{enumerate}

Tal curva é chamada a  \emph{braquistócrona.}
\end{problema}


\begin{problema}
Uma tira elástica fina é modelada, na aproximação de pequenas inclinações,
minimizando
\[
    E[y]=\int_0^1 y''(x)^2\,\d x
\]
entre funções suaves satisfazendo
\[
    y(0)=0,\quad y(1)=0,\quad y'(0)=\alpha,\quad y'(1)=\beta.
\]
Suponha que exista um minimizador. Mostre que ele tem de ser um polinômio cúbico, determinando-o explicitamente em termos
de $\alpha$ e $\beta$.
\end{problema}


%=====================================================================
\section*{12. Curvatura e geometria de curvas planas}
%=====================================================================

\noindent{\small A curvatura de uma curva mede a rapidez com que sua direção
tangente varia, e está governada por equações diferenciais simples. Os
problemas a seguir exploram esse ponto de vista.}

\vspace{4mm}

\begin{problema} Dada uma curva $\gamma:[0,L] \to \mathbb{R}^2,$ definimos sua \emph{curvatura com sinal} como a função 
\[
\kappa(s) = \frac{\det (\gamma'(s),\gamma''(s))}{\|\gamma'(s)\|^3},
\]
onde $(\gamma'(s),\gamma''(s))$ denota a matriz obtida fazendo $\gamma'(s) $ e $\gamma''(s)$ como seus vetores coluna. 
\begin{enumerate}[label=(\alph*)]
    \item Seja $f\in C^2$ e considere o gráfico $y=f(x)$. Mostre que sua
    curvatura com sinal é
    \[
        \kappa(x)=\frac{f''(x)}{\bigl(1+f'(x)^2\bigr)^{3/2}}.
    \]
    \item Calcule a curvatura de uma reta e de uma circunferência de raio
    $R$, e interprete geometricamente os resultados.
    \item (Espiral de Euler) Considere uma curva parametrizada por
    comprimento de arco cuja curvatura cresce linearmente, $\kappa(s)=s$.
    Resolva a equação intrínseca $\theta'(s)=\kappa(s)$ e expresse a curva
    $\gamma(s)=(x(s),y(s))$ por meio das integrais de Fresnel
    $\int_0^s\cos(u^2/2)\,\d u$ e $\int_0^s\sin(u^2/2)\,\d u$.
\end{enumerate}
\end{problema}

\begin{problema}
Seja $\kappa:[0,L]\to\R$ contínua. Defina
\[
    \theta(s)=\theta_0+\int_0^s \kappa(u)\,\d u,
\]
e em seguida a curva $\gamma(s)=(x(s),y(s))$ por
\[
    x'(s)=\cos\theta(s),
    \qquad
    y'(s)=\sin\theta(s),
    \qquad
    \gamma(0)=p_0.
\]
\begin{enumerate}[label=(\alph*)]
    \item Prove que $|\gamma'(s)|=1$, de modo que a curva está parametrizada
    por comprimento de arco.
    \item Definindo a curvatura com sinal por $\theta'(s)$, prove que a
    curvatura com sinal de $\gamma$ é $\kappa(s)$.
    \item Suponha que duas curvas parametrizadas por comprimento de arco
    tenham a mesma curvatura com sinal. Após uma rotação e uma translação de
    uma delas, suponha que ambas tenham o mesmo ponto inicial e o mesmo
    ângulo tangente inicial. Prove que as duas curvas coincidem.
\end{enumerate}
\end{problema}

\begin{problema}
Seja $\gamma:[0,L]\to\R^2$ uma curva $C^2$ parametrizada por comprimento de
arco, com curvatura com sinal constante $\kappa(s)=k$. Sejam
\[
    T(s)=\gamma'(s),
    \qquad
    N(s)=(-T_2(s),T_1(s)).
\]
\begin{enumerate}[label=(\alph*)]
    \item Mostre que $T'(s)=kN(s)$ e $N'(s)=-kT(s)$.
    \item Se $k=0$, prove que a imagem de $\gamma$ está contida em uma reta.
    \item Se $k\neq 0$, prove que $\gamma(s)+\frac{1}{k}N(s)$ é constante.
    Conclua que a imagem de $\gamma$ está contida em uma circunferência de
    raio $1/|k|$.
\end{enumerate}
\end{problema}

\begin{problema}
Seja $\gamma$ uma curva plana suave, simples, fechada e estritamente
convexa, parametrizada por comprimento de arco. Suponha que sua curvatura
com sinal $\kappa$ seja positiva em todos os pontos, e escreva
\[
    \gamma'(s)=(\cos\theta(s),\sin\theta(s)),
    \qquad
    \theta'(s)=\kappa(s).
\]
\begin{enumerate}[label=(\alph*)]
    \item Explique por que, ao percorrer a curva uma vez, a direção tangente
    dá exatamente uma volta completa; ou seja, $\theta(L)=\theta(0)+2\pi$
    para uma escolha contínua do ângulo tangente.
    \item Deduza que $\int_0^L \kappa(s)\,\d s=2\pi$.
    \item Explique brevemente por que a convexidade exclui voltas adicionais
    da direção tangente.
\end{enumerate}
\end{problema}

\begin{problema}
Sejam $a>0$ e $L>2a$. Entre todas as funções suaves
$y:[-a,a]\to[0,\infty)$ satisfazendo
\[
    y(-a)=y(a)=0,
    \qquad
    \int_{-a}^a \sqrt{1+y'(x)^2}\,\d x=L,
\]
considere o problema de maximizar a área $A[y]=\int_{-a}^a y(x)\,\d x$.
Suponha que exista um maximizador, positivo em $(-a,a)$. Mostre que ele tem
curvatura constante --- onde a curvatura com sinal de um gráfico é
\[
    \kappa=\frac{y''}{(1+y'^2)^{3/2}}
\]
--- e conclua que se trata de um arco de circunferência.
\end{problema}

\begin{problema}
Uma curva plana é parametrizada por comprimento de arco e tem curvatura
$\kappa(s)$. Seu ângulo tangente satisfaz
\[
    \theta'(s)=\kappa(s),\qquad
    x'(s)=\cos\theta(s),\qquad
    y'(s)=\sin\theta(s).
\]
Suponha que $\kappa$ seja contínua em $[0,\infty)$ e que
$\int_0^\infty |\kappa(s)|\,\d s<\infty$.
\begin{enumerate}[label=(\alph*)]
    \item Prove que $\theta(s)$ tem um limite finito $\theta_\infty$ quando
    $s\to\infty$.
    \item Explique por que o vetor tangente unitário $(x'(s),y'(s))$ tem
    limite, isto é, a curva possui uma direção assintótica.
    \item Suponha adicionalmente que
    $\int_0^\infty s|\kappa(s)|\,\d s<\infty$. Prove que
    \[
        \int_0^\infty |\theta(s)-\theta_\infty|\,\d s<\infty.
    \]
    \item Deduza que a curva permanece a uma distância limitada da reta que
    passa por $\gamma(0)$ na direção $(\cos\theta_\infty,\sin\theta_\infty)$.
\end{enumerate}
\end{problema}

\end{document}
